\phantomsection
\addcontentsline{toc}{chapter}{Conclusion}
\markboth{Conclusion}{Conclusion}
\chapter*{Conclusion}%

In this project thesis, the development and application of system-independent CVs using PCA were successfully demonstrated. 
By reducing the dimensionality of global atomic descriptors, a data-driven CV was constructed and applied to the saponification reaction. 
The results confirm that the reaction can be successfully replicated for this system using the PCA-based approach. 
The constructed CV effectively guided the MetaD simulation from the reactants to the products, capturing the essential transition states and complex mechanistic details, such as the water-assisted nucleophilic attack.

Despite these promising results, certain limitations were observed during the evaluation. 
The success of the driven simulations showed a noticeable dependence on the chosen MetaD parameters. 
While this parameter dependence is a normal characteristic of enhanced sampling techniques \cite{Laio2008}, 
it emphasizes that a comprehensive statistical analysis is required in future studies to determine the optimal simulation settings. 
Furthermore, the reconstructed FES derived from the bias potential was found to be inaccurate. 

Additionally, it was noted that the pre-trained MACE-MP-0 model is not the ideal MLIP for this specific use case.
Furthermore, it is important to remark that PCA is a linear dimensionality reduction method \cite{Jolliffe2002}. 
Since the true reaction path in the high dimensional descriptor space is highly likely to be non-linear, 
reducing it to a single dimension via one CV is a significant simplification.
Remarkably, despite these methodological simplifications and the suboptimal model choice, the approach remained effective and successfully drove the chemical reaction. 
This outcome could be an indication of the possible robustness and flexibility of the PCA-based CV construction. 
Overall, the proposed method provides a potentially highly capable, purely mathematical framework for reaction path analysis. 

